\subsection{CASIMIR ELEMENTS}

By Prop. 3 of §1, no. 3, there exist negative symmetric bilinear forms on g, separating and invariant under Ad(G) (if G is semi-simple, we can take the Killing form of g, for example). Let F be such a form. Recall (Chap. I, §3, no. 7) that the Casimir element associated to F is the element \(\Gamma\) of the centre of the enveloping algebra U(g) such that, for any basis \((e_{i})\) of g satisfying \(\mathrm{F}(e_{i}, e_{j}) = -\delta_{ij}\) , we have \(\Gamma = -\sum e_{i}^{2}\) .

In the remainder of this chapter we shall call the Casimir elements of G the elements of \(\mathrm{U}(\mathfrak{g})\) obtained in this way from negative separating invariant symmetric bilinear forms on g. If \(\Gamma\) is a Casimir element of G and if \(\tau: G \to \mathbf{GL}(V)\) is an irreducible representation of G, the endomorphism \(\Gamma_{V}\) of V is a homothety (Algebra, Chap. VIII, §3, no. 2, Th. 1), whose ratio we shall denote by \(\tilde{\Gamma}(\tau)\) .

PROPOSITION 4. Let \(\Gamma\) be the Casimir element of G.

\begin{enumerate}
\def\labelenumi{\alph{enumi})}
\item
  If \(\tau\) is an irreducible representation of \(\mathbf{G}\) , \(\tilde{\Gamma}(\tau)\) is real and positive. If \(\tau\) is not the trivial representation, \(\tilde{\Gamma}(\tau) > 0\) .
\item
  There exists a unique quadratic form \(\mathbf{Q}_{\Gamma}\) on \(\mathrm{X(T)}\otimes \mathbf{R}\) such that, for every irreducible representation \(\tau\) of \(\mathbf{G}\) ,
\end{enumerate}

\[
\tilde {\Gamma} (\tau) = \mathrm{Q} _ {\Gamma} (\lambda + \rho) - \mathrm{Q} _ {\Gamma} (\rho),
\]

where \(\lambda\) is the highest weight of \(\tau\) . The form \(Q_{\Gamma}\) is positive, separating and invariant under W.

Let \(\mathrm{F}\) be a separating negative symmetric bilinear form on \(\mathfrak{g}\) defining \(\Gamma\) . Let \(\tau : \mathrm{G} \to \mathbf{GL}(\mathrm{V})\) be an irreducible representation of \(\mathrm{G}\) ; let \(\langle , \rangle\) be a Hilbert scalar product on \(\mathrm{V}\) invariant under \(\mathrm{G}\) (§1, no. 1), and \((e_i)\) a basis of \(\mathfrak{g}\) such that \(\mathrm{F}(e_i, e_j) = -\delta_{ij}\) . Then, for every element \(v\) of \(\mathrm{V}\) not invariant under \(\mathrm{G}\) we have

\[
\begin{array}{c} \tilde {\Gamma} (\tau) \langle v, v \rangle = \langle v, \Gamma_ {\mathrm{V}} (v) \rangle = - \sum_ {i} \langle v, (e _ {i}) _ {\mathrm{V}} ^ {2} v \rangle = \sum_ {i} \langle v, (e _ {i}) _ {\mathrm{V}} ^ {*} (e _ {i}) _ {\mathrm{V}} v \rangle \\ = \sum_ {i} \langle (e _ {i}) _ {\mathrm{V}} v, (e _ {i}) _ {\mathrm{V}} v \rangle > 0, \end{array}
\]

hence \(a\) ).

Let B be the inverse form on \(t_{C}^{*}\) of the restriction to \(t_{C}\) of the bilinear form on \(g_{C}\) induced by F by extension of scalars. By the Cor. of Prop. 7 of Chap. VIII, §6, no. 4, we have \(\tilde{\Gamma}(\tau) = \mathrm{B}(\delta(\lambda), \delta(\lambda + 2\rho))\) . Extend \(\delta : \mathrm{X(T)} \to t_{\mathbf{C}}^{*}\) to an R-linear map from \(\mathrm{X(T)} \otimes \mathbf{R}\) to \(t_{C}^{*}\) and let \(Q_{\Gamma}\) be the quadratic form \(x \mapsto \mathrm{B}(\delta(x), \delta(x))\) on \(\mathrm{X(T)} \otimes \mathbf{R}\) ; it is separating and invariant under W, and

\[
\tilde {\Gamma} (\tau) = \mathrm{B} (\delta (\lambda + \rho), \delta (\lambda + \rho)) - \mathrm{B} (\delta (\rho), \delta (\rho)) = \mathrm{Q} _ {\Gamma} (\lambda + \rho) - \mathrm{Q} _ {\Gamma} (\rho).
\]

We show that the form \(Q_{\Gamma}\) is positive. Indeed, if \(x \in \mathrm{X}(\mathrm{T}) \otimes \mathbf{Q}\) , the element \(\delta(x)\) of \(t_{C}^{*}\) takes purely imaginary values on t, hence real values on it; we conclude by remarking that, for \(y \in it\) , we have \(\mathrm{F}(y, y) \geq 0\) .

It remains to prove the uniqueness assertion in b). Let Q be a quadratic form on \(X(T) \otimes \mathbf{R}\) satisfying the required condition, and let \(\Phi\) (resp. \(\Phi_{\Gamma}\) ) be the bilinear form associated to Q (resp. \(Q_{\Gamma}\) ). For \(\lambda, \mu \in X_{++}\) , we have

\[
\begin{array}{l} \Phi (\lambda , \mu) = (Q (\lambda + \mu + \rho) - Q (\rho)) - (Q (\lambda + \rho) - Q (\rho)) - (Q (\mu + \rho) - Q (\rho)) \\ \qquad = \Phi_ {\Gamma} (\lambda , \mu). \end{array}
\]

Since \(\mathrm{X(T)_{++}}\) generates the R-vector space \(\mathrm{X(T)}\otimes\mathbf{R}\) , we have \(\Phi=\Phi_{\Gamma}\) , so \(Q=Q_{\Gamma}\) .

Remark. Let \(x \in g\) . There exists a strictly positive real number A such that, for every irreducible representation \(\tau : G \to \mathbf{GL}(V)\) and every Hilbert structure on V invariant under G,

\[
\| \mathrm{L} (\tau) (x) \| ^ {2} \leq \mathrm{A}. \tilde {\Gamma} (\tau).
\]

Indeed, with the notations in the preceding proof, we can choose the basis \((e_{i})\) of g so that \(x = ae_{1}, a \in R\) . Then, for \(v \in V\) , we have

\[
\langle x _ {\mathrm{V}} v, x _ {\mathrm{V}} v \rangle = | a | ^ {2} \langle e _ {1} v, e _ {1} v \rangle \leq | a | ^ {2} \tilde {\Gamma} (\tau) \langle v, v \rangle .
\]

